\documentclass[10pt,a4paper]{article}
\usepackage[utf8]{inputenc}
\usepackage[T1]{fontenc}
\usepackage{lmodern}
\usepackage[margin=2.2cm]{geometry}
\usepackage{amsmath}
\usepackage{graphicx}
\usepackage{booktabs}
\usepackage{caption}
\usepackage{float}
\usepackage{hyperref}
\captionsetup{font=small,labelfont=bf}
\input{values}

\title{Few-electron charge states, and line fitting after the U-Net\\[2pt]
\large what was done, and what came out}
\author{\small Generated by \texttt{docs/charge\_states/make\_few.py} from
\texttt{scripts/run\_17\_few\_electron.py}\\
\small run: \texttt{\vRun}}
\date{}

\begin{document}
\maketitle

\section{What changed, and why}

The earlier results scored the whole $100\times100$ window, which holds 4 to
102 charge states per device.  An experiment cares about the few-electron
corner.  Two changes were made, and neither needs new training.

\begin{enumerate}
\item \textbf{Score only the few-electron region}: the pixels whose
  \emph{true} state has $n_1\le\vNmax$ and $n_2\le\vNmax$.  Labels are
  \emph{absolute}, with no shift, because the top-right corner is $(0,0)$ on
  every device and an experiment needs the real occupation.
\item \textbf{Fit lines after the U-Net} (\texttt{fit\_lines}).  The U-Net
  maps fail through gaps: one gap joins two states into one region, and every
  label behind it shifts.  The line fit repairs this:
  \begin{itemize}
  \item thin the binary map and cut it at its junctions into pieces;
  \item replace every piece by its least-squares straight segment;
  \item extend each \emph{free} end (one with nothing attached) along the
    segment's own direction, and draw the extension only if it reaches
    another line or the frame within a set length.  A line is joined to what
    it was heading for and never left dangling.
  \end{itemize}
  The knobs are the shortest piece kept, the longest extension, and whether
  the original map is kept underneath.  All \vNFit{} settings were scanned on
  the 75 validation devices, over every U-Net, and the best by balanced
  accuracy was applied unchanged to the test devices: \emph{\vFitKnob}.  On
  validation it gives \vFitVal\,\% balanced accuracy, against \vPlainVal\,\%
  without the fit.  The chosen setting is at the edge of the scanned range
  (shortest piece and longest extension), so a wider scan might do slightly
  better.
\end{enumerate}

\subsection*{Metrics, in the few-electron region}
\begin{itemize}
\item \textbf{pixel accuracy}: the fraction of the region's scored pixels
  with the right $(n_1,n_2)$.  It is dominated by the large $(0,0)$ area.
\item \textbf{balanced accuracy}: the same accuracy computed per true state,
  then averaged.  Every state counts once, however large it is.  This is the
  headline number.
\item \textbf{(1,1) landing}: park the gates at the centre of the
  \emph{predicted} $(1,1)$ region (its point farthest from the region's
  edge).  Is the device really in $(1,1)$?  It scores 0 if no $(1,1)$ was
  predicted.  It is taken over the \vNLand{} of \vNDev{} test devices whose
  window contains $(1,1)$.
\end{itemize}
Pixels within 1\,px of a true line are not scored.

\subsection*{Methods compared}
\begin{itemize}
\item \textbf{all (0,0)}: label everything $(0,0)$; the floor.
\item \textbf{dense, no ML}: every pixel of the sensor image measured;
  lines found by the step-residual detector (chosen on validation in the
  earlier study).
\item \textbf{full CSD}: the true line map; this only checks the labelling.
\item \textbf{U-Net} and \textbf{U-Net + line fit}, for each of the 15
  measurement budgets.
\end{itemize}

\section{Results (\vNDev{} test devices)}

\begin{figure}[H]
\centering
\includegraphics[width=\textwidth]{figures/fig_few_vs_coverage.pdf}
\caption{Few-electron region ($n_1,n_2\le\vNmax$).  Each point is the mean
over the test devices for one U-Net budget.  Blue: U-Net map as is.  Green:
after the line fit.  Orange: dense sensor image with no network.  Star: true
line map.  Dashed: labelling everything $(0,0)$.}
\label{fig:cov}
\end{figure}

\begin{table}[H]
\centering
\caption{Test split.  Each cell: mean [median] over devices.  Few-electron
columns use absolute labels in the region $n_1,n_2\le\vNmax$.  The
whole-window columns are the earlier relative accuracy and state recall.}
\label{tab:main}
\resizebox{\textwidth}{!}{\input{summary.tex}}
\end{table}

\subsection*{Main findings}
\begin{itemize}
\item \textbf{The line fit roughly doubles what the U-Net maps give.}  At
  \vBud{} (\vCov\,\% of the grid) the balanced accuracy rises from \vUBal\,\%
  to \vFBal\,\%, the pixel accuracy from \vUAcc\,\% to \vFAcc\,\%, and (1,1)
  landing from \vULand\,\% to \vFLand\,\%.  The balanced accuracy improves for
  \vGainPos{} of \vNBud{} budgets, by \vGainMin{} to \vGainMax{} points;
  landing improves by \vLGainMin{} to \vLGainMax{} points.
\item \textbf{The best combination}, \vBestF{} with the line fit, reaches
  \vBestFBal\,\% balanced accuracy and \vBestFLand\,\% landing.  The best
  U-Net without the fit reaches \vBestUBal\,\% (\vBestU).
\item \textbf{The dense measurement with no network is still far ahead}:
  \vDenseBal\,\% balanced accuracy and \vDenseLand\,\% landing, but it
  measures every pixel; the U-Nets measure \vLoCov--\vHiCov\,\%.
\item \textbf{Landing is the hard part.}  Even after the fit, the gates land
  in $(1,1)$ on fewer than half of the devices for every budget, and the
  median landing is 0 for every U-Net row.  For tuning to $(1,1)$ the
  U-Net maps are not yet reliable.
\item \textbf{The floor}: ``all (0,0)'' gives \vBaseAcc\,\% pixel accuracy in
  this region but only \vBaseBal\,\% balanced accuracy and \vBaseLand\,\%
  landing.  This is why balanced accuracy and landing are the numbers to
  quote, not pixel accuracy.
\end{itemize}

\begin{figure}[H]
\centering
\includegraphics[width=\textwidth]{figures/fig_fit_examples.pdf}
\caption{Three test devices, \vBud{} U-Net.  (a) U-Net line map.
(b) Labels from (a).  (c) After the line fit.  (d) Labels from (c).
(e) Truth, with the few-electron region outlined (on sample\_223 the whole
window is inside it).  Top: the fit closes the gaps and recovers most of the
states; one extension crosses the $(1,0)$ region, so the counting there is
off by one.  Middle: a complete map is left unchanged.  Bottom: dashed
dot-1 lines are too broken for straight segments to rejoin, and the fit does
not help.  Negative labels such as $-1$ mark where the counting went
wrong.}
\label{fig:ex}
\end{figure}

\section{Caveats}
\begin{itemize}
\item \textbf{One sweep, noise-free images.}  Only the 4--8-ray sweep is
  available.  With experimental noise both the dense detector and the U-Net
  would do worse.
\item \textbf{The fit can draw false lines.}  An extension that misses its
  target and reaches a different line splits a region in two
  (Fig.~\ref{fig:ex}, top).  Keeping it within a set length limits this but
  does not prevent it.
\item \textbf{The corner is assumed to be $(0,0)$.}  This holds for every
  device in this simulated pool.
\item \textbf{$\le\vNmax$ is a choice.}  It is one setting (\texttt{NMAX})
  in \texttt{run\_17\_few\_electron.py}.
\end{itemize}

\end{document}
